% Generated draft: requires mathematical and source review.
\documentclass[11pt]{article}
\usepackage[margin=25mm]{geometry}
\usepackage{fontspec}
\setmainfont{Latin Modern Roman}
\usepackage{amsmath,amssymb,mathtools,unicode-math}
\setmathfont{Latin Modern Math}
\usepackage{enumitem,needspace,xurl,hyperref}
\setlength{\parindent}{0pt}
\setlength{\parskip}{5pt}
\setlength{\emergencystretch}{3em}
\newcommand{\N}{\mathbb{N}}
\newcommand{\Z}{\mathbb{Z}}
\newcommand{\Q}{\mathbb{Q}}
\newcommand{\R}{\mathbb{R}}
\newcommand{\C}{\mathbb{C}}
\newcommand{\F}{\mathbb{F}}
\newcommand{\A}{\mathbb{A}}
\newcommand{\PP}{\mathbb{P}}
\newcommand{\E}{\mathbb{E}}
\newcommand{\eps}{\varepsilon}
\newcommand{\dd}{\,\mathrm{d}}
\begin{document}
\section*{UnsolvedMath ID: 30005286}
\textbf{Source problem code:} catalogue entry
\subsection{Definitions and mathematical statement}
The source uses \(\operatorname{SDCF}\) for separably differentially closed fields and refers to separably closed fields, abbreviated \(\operatorname{SCF}\), as a guide to further model-theoretic questions. It also names a theory rendered in the supplied cleaned statement as \(\operatorname{SDCF}_{\ell,p,\epsilon}\). The record does not specify its axioms, language, the meanings or permitted values of \(\ell,p,\epsilon\), or the derivation structure. Consequently, these symbols retain their source-dependent meanings here; no characteristic, number of derivations, or parameter range is imposed.

For clarity about the word ``types,'' if \(T\) is a theory in a language \(L\), \(M\models T\), \(A\subseteq M\), and \(n\geq 1\), a complete \(n\)-type over \(A\) is a maximal consistent collection of \(L(A)\)-formulas in \(n\) free variables, consistent with the complete theory of \((M,a)_{a\in A}\). Write \(S_n(A)\) for the set of these types. This standard definition supplies terminology only: the source does not specify whether its question concerns complete types, partial types, particular arities, or particular parameter sets.

\textbf{Source research questions.} With the intended theories and differential-algebraic conventions fixed, address the following directions:
\begin{enumerate}
\item Determine further model-theoretic properties of \(\operatorname{SDCF}\), taking the theory of separably closed fields as a guide. The source gives no exhaustive list of properties to establish.
\item Describe types of \(\operatorname{SDCF}_{\ell,p,\epsilon}\) in terms of separable ideals in a suitable differential polynomial ring. The coefficient field, differential polynomial ring, meaning of ``separable ideal,'' relevant class of ideals, and intended correspondence with types must be specified before this becomes a precise classification assertion.
\item Determine whether there is a form of Zilber's dichotomy in this context, and formulate and establish such a dichotomy if appropriate. The record does not specify the objects to which it should apply, its hypotheses, or its alternatives.
\end{enumerate}
These are open-ended research questions in the supplied record, rather than a fully specified quantified conjecture. In particular, this draft does not assert a correspondence between types and ideals or select a particular formulation of Zilber's dichotomy.
\subsection{Short English statement}
Investigate the model theory of separably differentially closed fields, seek descriptions of types using separable differential ideals, and ask whether a suitable Zilber dichotomy holds. The source leaves the theory parameters, ideal conventions, scope of types, and proposed dichotomy unspecified, so a precise theorem statement requires further source clarification and human mathematical review.
\subsection{Sources}
\begin{itemize}
\item Local UnsolvedMath registry, permanent ID \(30005286\), title \emph{Model Theory and Zilber Dichotomies for SDCF}, source problem code \texttt{catalogue entry}, in the Oberwolfach Reports Open Problems collection. No individual proposer is identified in the supplied record.
\item Supplied report citation: \emph{Mini-Workshop: Topological and Differential Expansions of o-minimal Structures} (2023), Oberwolfach Reports, pp.~3005--3050, DOI \texttt{10.4171/owr/2022/52}. \url{https://doi.org/10.4171/owr/2022/52}. The original statement includes the page marker 3033. The registry assigns proposed year 2022; the relationship between that year and the citation's 2023 date is not explained in the record. The report is described as documenting workshop discussions and open problems.
\item Supplied literature reference: Ino--Sánchez, \emph{Separably differentially closed fields}. \url{https://arxiv.org/abs/2302.11319}.
\item Supplied literature reference: EMS Press, OWR 2022/52. \url{https://doi.org/10.4171/owr/2022/52}.
\item Literature assessment dated exactly \texttt{2026-08-22}: the supplied curation labels these directions open and reports that existing work gives completeness and quantifier-elimination results, but not the requested full classification. This is supplied dated curation, described as literature triage derived from the supplied OWR curation, and has not been independently verified here. Its time-sensitive status and cited claims require independent verification.
\item The registry labels the statement \texttt{corrected\_verified} and reports that it was reconstructed and checked against the official Oberwolfach report. These are supplied provenance claims; this draft does not independently verify them or claim to have visited either reference.
\end{itemize}
\end{document}
